% Chapter 30: The Second and Third Laws of Thermodynamics (complete text)
\chapter{The Second and Third Laws of Thermodynamics}

\step{A Counterintuitive Opening}

Begin with an absurdly large energy account.

Cooling a kilogram of seawater by one degree releases about 4200 joules. Earth's oceans contain roughly a hundred septillion, \(1.4\times10^{21}\), kilograms of water. Cool them all by one degree and the released heat is about \(6\times10^{24}\) joules, compared with humanity's annual consumption of roughly \(6\times10^{20}\) joules. That one-degree cooling could supply us at present rates for about ten thousand years.

Imagine a ship needing no coal, oil, or nuclear fuel. As it sails, it extracts seawater heat, generates electricity for its propeller, and discharges slightly cooler water. Losing a little heat scarcely affects the ocean; another ship follows. This ship \textbf{does not violate energy conservation}: every unit of work is paid for by reduced seawater internal energy. The ledger balances perfectly.

Has anyone built it? Not one. More strongly, not even one screw operates this way. The failure is peculiar: neither low efficiency nor high cost, but impossibility in principle, unable to complete even its first turn. Nineteenth-century engineers tried many schemes, all wrecked on the same reef.

Such a machine is a \textbf{perpetual-motion machine of the second kind}. It is subtler than the first kind, which creates energy from nothing. The first cheats openly with unbalanced books; the second keeps impeccable accounts yet remains impossible. Catching it requires another legal code.

The reef lies in laws, not engineering. Energy conservation, Chapter 28's subject, governs quantity but not direction. Two directional laws now enter: the second explains heat's one-way use and the iron ceiling on engine efficiency; the third explains why we can never stand on temperature's floor.

\step{Make a Guess First}

Before reading on, guess why the seawater-powered ship is impossible.

First: technology is inadequate. As ancient people could not build aircraft, perhaps future engineers can build this ship.

Second: seawater has abundant but dilute heat. Low-temperature heat is hard to use; concentrate it into high-temperature heat and the ship can sail.

Third: a new law independent of energy conservation forbids producing work solely from one heat reservoir, beyond all technological escape.

Also guess a household question: would opening a refrigerator door cool a hot summer kitchen, warm it, or leave it unchanged? Record both judgments for checking at the end.

\step{Hands-on Experiment}

\begin{experiment}[Can an Open Refrigerator Be an Air Conditioner?]
% BEGIN TRANSLATOR SAFETY NOTE
\textbf{Translator safety note:} Treat the open-door test as a thought experiment: do not tape down sensors, defeat controls, or leave a food-filled refrigerator open for hours. Keep perishable food properly chilled. See \href{https://www.fda.gov/food/food-safety-during-emergencies/power-outages-key-tips-consumers-about-food-safety}{FDA refrigeration and food-safety guidance}.
% END TRANSLATOR SAFETY NOTE

You need a working refrigerator, a thermometer capable of room-temperature measurement, not a clinical thermometer, and a phone timer.

First place the thermometer more than two meters from the refrigerator, wait five minutes, and record room temperature.

Second leave the refrigerator door open. To keep it cooling, tape down the little door-closed sensor switch on the frame if there is one, keeping the compressor running. Wait one or two hours without operating other high-power appliances.

Third read the thermometer again. Also touch the refrigerator's back or sides: hot or cool? Check whether the interior really is cooling.

You will see that the interior does cool, but mean room temperature does not fall; it rises slightly. After an hour or two of hard work, the refrigerator has warmed the kitchen.
\end{experiment}

Pause for three seconds. Why does a cooling refrigerator warm the room? It extracts heat from the cold interior and rejects it through rear coils into the warmer room. Heat goes cold to hot, apparently backward. But the compressor continuously consumes electricity, which ultimately also becomes heat in the same room. The room receives transported interior heat \textbf{plus} electrically purchased heat. Both are positive, so it warms.

Why does real air-conditioning cool a room? Its outdoor heat-rejection unit discharges outside. Cooling requires an outside to receive heat. This simple observation, \textbf{moving heat cold to hot requires extra work, with a larger environment receiving the bill}, is the second law's everyday face.

Record answers to three observations for later. Does the back's heat come from transported interior heat or compressor production? In a perfectly insulated room, would warming eventually stop if the refrigerator stayed open? Its interior cools while the whole warms: how does this relate to direction?

\step{The Old Model Runs into Trouble}

Return to Chapter 28. The first law balances internal-energy change against heat and work: energy neither appears nor disappears. It vetoes every energy-from-nothing machine, the deservedly doomed first kind of perpetual motion.

Yet ask it to judge three processes and it cannot reject any:

Cool coffee gathers room heat and boils again: the room loses exactly what the coffee gains. Ice in an ice--water mixture gets colder while surrounding water gets hotter: water's lost heat exactly equals ice's gained heat, merely with direction reversed. The seawater ship also balances, as already shown.

All conserve energy; none has ever been seen. The first law is an accountant checking totals without examining flow direction. A new law must answer: \textbf{among all energy-conserving processes, which really occur?}

Interestingly, its first foundation preceded energy conservation. In 1824, twenty-eight-year-old French engineer Sadi Carnot published a booklet asking whether steam-engine efficiency had an upper bound. Engines achieved only a few percent, and engineers did not know whether the missing ninety-plus percent awaited improvement or was fundamentally unavailable. An idealized thought experiment answered: \textbf{there is an upper bound, depending only on the two operating temperatures, not working substance or structure}. Ironically Carnot still believed caloric theory, deriving a correct conclusion within a wrong framework. Good thought experiments can outlive bad theories. Over twenty years later, Clausius and Kelvin rebuilt the insight on energy conservation, giving the second law.

\begin{mainline}
The difficulty in one sentence: \textbf{energy conservation governs total balance, not each transaction's direction}. The three legal-but-unseen processes share one directional issue: heat spontaneously flows hot to cold, never backward. The new law must address direction.
\end{mainline}

\step{Inventing New Concepts}

History leaves two classic statements of the second law, faces of one coin.

\textbf{Clausius's statement}: heat cannot transfer from a colder to a hotter object without other changes. Notice that final qualification. Refrigerators do move heat uphill, but cause other changes: the meter runs and electricity becomes heat. The forbidden reverse flow is \textbf{cost-free}.

\textbf{Kelvin's statement}: heat cannot be taken from a single reservoir and wholly converted to work without other effects. Single reservoir strikes directly at our ship: seawater is its only reservoir, with no colder place for waste heat, so none of its absorbed heat becomes net work. It does not forbid heat becoming work, which engines do daily. It forbids \textbf{complete conversion}, and conversion \textbf{from only one reservoir}.

These statements are equivalent. Suppose a machine violates Kelvin, wholly converting one reservoir's heat to work. Use that work to drive an ordinary refrigerator, moving cold-store heat into that reservoir. The net effect is cost-free heat flow cold to hot, violating Clausius too. Conversely, a Clausius-violating machine plus an ordinary engine violates Kelvin. The statements stand or fall together.

What about Carnot's ceiling? His perfect imaginary engine contacts two constant-temperature reservoirs, hot at \(T_1\) and cold at \(T_2\). It first absorbs hot-reservoir heat, then disconnects and expands, doing work and cooling; next it rejects heat to the cold reservoir and finally compresses back to its starting state. Every stage is so slow as to remain nearly equilibrated, without friction or temperature-difference turbulence, reversible at any moment. This \textbf{Carnot cycle} defines a reversible engine. Carnot proved that all reversible engines between the same reservoirs have equal efficiency, and no real irreversible engine can exceed it. Depending only on the two temperatures, that efficiency even supplies a substance-independent thermometer: comparing two engine efficiencies defines temperature ratios. This thought-thermometer's scale is the absolute scale, in kelvins.

\begin{center}
\begin{tikzpicture}[scale=1.0, >={Stealth[length=2.5mm]}]
  % Hot reservoir
  \draw[thick, fill=red!12] (-2.2,3.4) rectangle (2.2,4.2);
  \node at (0,3.8) {Hot reservoir \(T_1\)};
  % Cold reservoir
  \draw[thick, fill=blue!10] (-2.2,-1.0) rectangle (2.2,-0.2);
  \node at (0,-0.6) {Cold reservoir \(T_2\)};
  % Engine
  \draw[thick, fill=yellow!15] (0,1.6) circle (0.85);
  \node at (0,1.6) {Engine};
  % Q1 arrow
  \draw[->, very thick, red!60!black] (-0.55,3.4) -- (-0.55,2.5);
  \node[left] at (-0.6,2.95) {Heat in \(Q_1\)};
  % Q2 arrow
  \draw[->, very thick, blue!60!black] (0.55,0.7) -- (0.55,-0.2);
  \node[right] at (0.62,0.25) {Heat out \(Q_2\)};
  % W arrow
  \draw[->, very thick] (0.85,1.6) -- (2.1,1.6);
  \node[right] at (2.15,1.6) {Work out \(W=Q_1-Q_2\)};
\end{tikzpicture}
\end{center}

Study this central picture for ten seconds: an engine is \textbf{not merely a heat-to-work converter, but a waterwheel intercepting some heat as it flows hot to cold}. How much becomes work, \(W\), and how much leaves as waste heat, \(Q_2\), depends on two temperatures. Waste heat is not poor engineering but necessary operation: without downstream, no flow.

Chapter 29 gives a deeper statement: isolated-system entropy almost always rises and reaches a maximum at equilibrium, \(\Delta S \geq 0\). Why does hot-to-cold transfer raise entropy? The same heat \(Q\) accounts for less per degree at high temperature and more at low temperature, so the transfer profits the total account. Next we calculate it.

\step{The Model in One Sentence}

\begin{oneline}
An engine intercepts part of heat flowing hot to cold as work. The captured fraction has a strict ceiling set only by the two temperatures. Thus a single reservoir's heat cannot yield even a penny of work, and approaching absolute zero makes remaining heat ever harder to reject: zero can be approached but never reached.
\end{oneline}

Three forevers hide here: some waste heat always leaves, Kelvin; reverse flow always costs, Clausius; zero always remains another step away, the third law. These are not three bans but three faces of Chapter 29's entropy statistics, appearing in engines, refrigerators, and thermometers.

\step{The Mathematical Translator}

Give thought numbers. A reversible engine draws from hot reservoir \(T_1\) heat \(Q_1\), and rejects to cold reservoir \(T_2\) heat \(Q_2\), and outputs \(W=Q_1-Q_2\). Efficiency is received work divided by supplied heat:
\[
\eta=\frac{W}{Q_1}=\frac{Q_1-Q_2}{Q_1}=1-\frac{Q_2}{Q_1}.
\]

Reversibility produces no entropy over the cycle. The hot reservoir gives heat \(Q_1\), losing entropy \(Q_1/T_1\); the cold receives \(Q_2\), gaining \(Q_2/T_2\). Exact cancellation requires
\[
\frac{Q_1}{T_1}=\frac{Q_2}{T_2}
\quad\Longrightarrow\quad
\frac{Q_2}{Q_1}=\frac{T_2}{T_1}.
\]
Substitute for the Carnot limit:
\[
\boxed{\ \eta \leq 1-\frac{T_2}{T_1}\ }
\]
Temperatures must be absolute, in kelvins measured from absolute zero. Celsius's arbitrary freezing-point zero cannot enter this formula.

Check special cases immediately. First, \(T_2=T_1\): no difference gives \(\eta=0\). Equal-temperature reservoirs cannot run an engine, as a level river cannot power a hydroelectric station. The seawater ship sinks again. Second, \(T_2\to 0\) gives \(\eta\to 100\%\), but the third law says a \(T_2=0\) reservoir does not exist, keeping efficiency just below perfection. The laws interlock. Third, reverse direction: \(T_2>T_1\) gives negative efficiency. Absurd? No: the machine runs backward, taking work to move heat cold to hot. An engine becomes a refrigerator. The formula's reversal matches the machine's, checking its form.

Below is the Carnot cycle in the pressure--volume plane. Chapter 28 showed that enclosed area equals net cycle work:

\begin{center}
\begin{tikzpicture}[scale=1.15, >={Stealth[length=2.2mm]}]
  \draw[->] (0,0) -- (5.6,0) node[right] {Volume \(V\)};
  \draw[->] (0,0) -- (0,3.9) node[above] {Pressure \(p\)};
  % Isothermal heat input T1: A->B
  \draw[thick, red!60!black] plot[smooth, tension=0.9] coordinates {(1.0,3.3) (1.6,2.7) (2.4,2.15) (3.2,1.75)};
  % Adiabatic expansion: B->C
  \draw[thick] plot[smooth, tension=0.9] coordinates {(3.2,1.75) (3.9,1.15) (4.6,0.72)};
  % Isothermal heat rejection T2: C->D
  \draw[thick, blue!60!black] plot[smooth, tension=0.9] coordinates {(4.6,0.72) (3.2,0.85) (1.7,1.15)};
  % Adiabatic compression: D->A
  \draw[thick] plot[smooth, tension=0.9] coordinates {(1.7,1.15) (1.3,2.1) (1.0,3.3)};
  \fill (1.0,3.3) circle (0.05) node[above left] {A};
  \fill (3.2,1.75) circle (0.05) node[above right] {B};
  \fill (4.6,0.72) circle (0.05) node[below right] {C};
  \fill (1.7,1.15) circle (0.05) node[below left] {D};
  \node[red!60!black] at (3.2,3.65) {\shortstack{Isothermal heat in\\\(Q_1\) (\(T_1\))}};
  \node[blue!60!black] at (3.2,-0.75) {\shortstack{Isothermal heat out\\\(Q_2\) (\(T_2\))}};
  \node[rotate=-32] at (4.65,1.95) {\small Adiabatic expansion};
  \node[rotate=62] at (0.45,2.1) {\small Adiabatic compression};
\end{tikzpicture}
\end{center}

A to B: expand slowly against the hot reservoir, absorbing \(Q_1\). B to C: disconnect and expand until temperature falls to \(T_2\). C to D: compress against the cold reservoir, releasing \(Q_2\). D to A: disconnect and compress back to \(T_1\), closing the cycle. B--C and D--A exchange no heat: the adiabatic stages. The machine communicates externally at only two temperatures, the secret of reversibility: every step can retrace its path.

\begin{mathbridge}
We jumped from reversibility to Carnot efficiency through unchanged entropy, but ideal gas provides a stepwise derivation showing how temperature enters. Isothermal internal energy is constant, so heat becomes expansion work: \(Q_1=nRT_1\ln(V_B/V_A)\), and similarly \(Q_2=nRT_2\ln(V_C/V_D)\). The adiabatic stages satisfy \(T V^{\gamma-1}=\text{constant}\), giving \(V_B/V_A=V_C/V_D\) on comparison. Equal logarithms yield \(Q_2/Q_1=T_2/T_1\). Every gas property, \(n\), \(\gamma\), and volume, cancels, leaving temperature ratio: algebraic evidence of working-substance independence.

A refrigerator has a corresponding ceiling. Reverse Carnot, supply work \(W\), and use cold reservoir \(T_2\) to extract heat \(Q_2\). Its coefficient of performance is
\[
\mathrm{COP}_{\text{cooling}}=\frac{Q_2}{W}\leq\frac{T_2}{T_1-T_2}.
\]
Check absolute zero. Fix room temperature \(T_1\) and lower the cold reservoir's \(T_2\). COP tends to zero: removing the same heat from colder places costs more work, approaching infinite work for even a little heat near zero. \textbf{The third law's unattainable zero is already written in this formula.}
\end{mathbridge}

Now two real-number estimates.

First, a fossil-fuel power station. Modern ultra-supercritical steam is about \(600\,^\circ\mathrm{C}\), \(873\ \mathrm{K}\), while condensers cooled by river water or towers operate near \(30\,^\circ\mathrm{C}\), \(303\ \mathrm{K}\). Carnot gives
\[
\eta\leq 1-\frac{303}{873}\approx 65\%.
\]
The best actual plants achieve about \(45\%\), losing portions to material-temperature limits, transfer temperature differences, and friction. Carnot promised a ceiling, not a guarantee. Use it backward for power: someone claiming seventy-percent efficiency with \(600\,^\circ\mathrm{C}\) steam and \(30\,^\circ\mathrm{C}\) surroundings violates the second law, no blueprints needed.

Second, geothermal power. Water emerges at \(150\,^\circ\mathrm{C}\), \(423\ \mathrm{K}\), into surroundings at \(25\,^\circ\mathrm{C}\), \(298\ \mathrm{K}\). The limit is \(1-298/423\approx 30\%\), while actual plants often achieve only ten-odd percent. Low-temperature sources inherit low ceilings, explaining why attractive-sounding waste-heat generation often disappoints economically.

Finally introduce temperature's floor. The \textbf{third law of thermodynamics} has two equivalent statements. Nernst's form says a perfect crystal's entropy approaches zero as temperature approaches absolute zero. In Chapter 29's language, one microstate remains, \(\Omega\to1\), with no larger room available. Unattainability says no finite sequence of operations can cool any system to absolute zero. Intuitively, cooling rejects heat to something colder, but the coldest has no colder recipient. Near zero, less heat remains removable, transfer slows, and COP shrinks as above. Each step grows harder: infinite stairs, finite steps. Experimental history is continual approach: liquid helium gives \(4.2\ \mathrm{K}\), dilution refrigeration about \(0.01\ \mathrm{K}\), and laser-cooled atomic clouds reach order a billionth of a kelvin. Records advance; zero forever remains one final step away.

\step{The Model's Limits}

As usual: assumptions, domain, and misuses.

\textbf{The second law is statistical, not absolute.} Chapter 29 showed entropy increase as near-certainty, not prohibition. Small systems permit brief, local reversals: molecular impacts can momentarily give micrometer suspended particles work from thermal motion. This \textbf{does not violate} a law about macroscopic averages. But do not get excited about harvesting fluctuations for a steady nano-engine: require \textbf{continuous, reliable} output and the average returns to entropy increase. Nearly impossible versus absolutely forbidden is its deepest character: it governs large-number worlds.

\textbf{Carnot is a criterion, not an engineering promise.} It assumes only two fixed-temperature reservoirs, infinitely slow quasistatic operation, no friction or leakage, and full reversibility. Real power output needs finite temperature differences for heat transfer, itself producing entropy. Speed and economy conflict. Use the ceiling to expose impossible claims and guide improvement where actual efficiency remains far below it.

\textbf{Do not muddle heat and temperature in slogans.} Converting all environmental heat to work fails not from insufficient heat but from no colder entropy sink. Conversely, smart winter heating uses electricity not merely to make heat, one unit for one, but to \textbf{move} several heat units with one electrical unit. Three outputs from one input may alarm intuition, yet extra heat comes outdoors, with electricity paying transport. This is the second law's gift, not a loophole.

\begin{challenge}
A subtler exceptional domain is self-gravitation. Drag makes satellites lose energy and descend while speeding up; contracting star clusters warm. Gravitating systems have negative heat capacity, requiring careful reformulation of ordinary second-law statements. More remarkably, black holes have entropy proportional to area and a temperature, radiating extremely slowly. Thermodynamics' three laws apply almost word for word, with gravity added to the ledger. This thread reaches unresolved general-relativistic and quantum-gravity territory; we return later.
\end{challenge}

\begin{mainline}
Instructions: in macroscopic, near-equilibrium settings with clear reservoirs, use second-law direction and Carnot bounds confidently. Small systems need fluctuations; strong gravity needs revised rules. Confusing the Carnot ceiling with actual efficiency, statistics with an absolute ban, or a heat pump's over-one-hundred-percent performance with perpetual motion are the three classic misuses.
\end{mainline}

\step{Explaining the Opening Phenomenon}

Now judge the seawater ship.

Its only reservoir is seawater. Carnot requires a colder place to receive rejected heat for net work output. Is there one? No: nearby water, air, and ship are at roughly one temperature. Without downstream, the waterwheel cannot even idle. Kelvin forbids this complete single-reservoir conversion. Entropy says it more sharply: seawater cooling loses entropy, zero rejected heat increases nobody else's entropy, and the isolated total decreases. Probability condemns it, not merely as difficult but as Sun-reaching-decimal-zeros impossible (Chapter 29).

One further blow: if it existed, connect its output to an ordinary refrigerator and heat could flow automatically from cold seawater to warmer air. Clausius falls too, cool coffee boils, and air-conditioning needs no power. One ship's success would sink the entire law's fleet. If the law stands, the ship never exists. This is direction, not technology.

The open-fridge answer is equally clear. It already moves heat uphill from cold interior to warmer room, paying electricity, which also remains as room heat. Transported heat plus electricity-derived heat are both positive, so the room warms. Its hot back displays Clausius's receipt. Real air-conditioning cools because an outside exists, where its outdoor unit rejects heat to atmosphere.

Check the three guesses. Inadequate technology is wrong: no precision makes a law step aside, as no perfect lever defeats momentum conservation. Concentrating low-temperature heat is wrong because concentrating it means pumping cold to hot, paying more work than earned; Carnot ensures a losing bargain. The third is right: the second law is independent of energy conservation, governing direction rather than total quantity.

\step{Thought Experiments and Challenges}

\begin{thought}[The Sun as Boiler, the Galaxy as Condenser]
Design a cosmic engine. Use the Sun's surface at about \(5800\ \mathrm{K}\) as boiler and the universe as condenser. Deep space contains the microwave background at only \(2.7\ \mathrm{K}\), a great cold sink for radiation. Carnot gives
\[
\eta\leq 1-\frac{2.7}{5800}\approx 99.95\%.
\]
Nearly perfect: the ultimate ceiling for starlight power. But one final obstacle remains. Even your condenser is at \(2.7\ \mathrm{K}\), so waste heat must be rejected above \(2.7\ \mathrm{K}\). The last \(0.05\%\) can never be recovered. Expansion keeps cooling the background, yet it never reaches zero. The engine's efficiency ceiling and refrigerator's temperature floor meet under one sky.
\end{thought}

\begin{thought}[Do Refrigerators and Life Violate the Second Law?]
Ice forms orderly crystals in refrigerators; fertilized eggs become intricately structured people. Local entropy drops greatly. Is the law wrong? You have the tools: inspect boundaries. The refrigerator interior is not isolated and rejects more heat into the kitchen. A person is not isolated; eating, breathing, and heat loss ensure total person-plus-environment entropy grows. Local order is purchased with greater surrounding disorder. Chapter 32 expands this large theme, with Maxwell's demon collecting the bill.
\end{thought}

\begin{puzzles}
\item \textbf{A heat pump over one hundred percent.} An advertisement promises four thermal kilowatt-hours indoors per electrical kilowatt-hour in winter, \(400\%\) heating performance, versus a resistance heater's maximum \(100\%\). Does this violate conservation? Why does performance fall markedly in severe cold, below minus ten degrees Celsius?

  \textit{Hint:} No violation. Resistance heaters convert electricity one-for-one; heat pumps transport outdoor heat, with the extra three units coming from outside. The Carnot heating coefficient is bounded by \(T_1/(T_1-T_2)\). A larger indoor--outdoor difference increases the denominator and transport difficulty. The formula already contains the weather's influence.

\item \textbf{A cup's temperature-difference plot.} Partition water at \(50\,^\circ\mathrm{C}\), somehow cool the left half to \(40\,^\circ\mathrm{C}\) and warm the right to \(60\,^\circ\mathrm{C}\), then run an engine from the difference until both return to \(50\,^\circ\mathrm{C}\), repeating forever. Energy seems conserved. Which step is fatal?

  \textit{Hint:} Not the engine, but making one half colder and the other hotter: a cost-free temperature difference reverses spontaneous heat flow and violates Clausius. The innocent engine never receives its upstream supply. Judge each step's direction, not merely total balance.

\item \textbf{A relay to zero.} One refrigerator cannot reach absolute zero, but what about a relay: A reaches \(1\ \mathrm{K}\), B reaches \(10^{-3}\ \mathrm{K}\), C goes lower, and so on? Why cannot successive closer steps arrive?

  \textit{Hint:} Each cooling stage needs a still colder reservoir to accept waste heat, precisely what you are trying to create. Meanwhile COP scales as \(T_2/(T_1-T_2)\), so removing fixed heat costs approaching infinite work near zero. Infinite stairs cannot be traversed in finitely many real operations, the third law's unattainability statement. Approach is allowed; arrival is not.

\item \textbf{Exhaust-recovery limits.} Someone proposes recovering all car-exhaust heat as power to reach one-hundred-percent efficiency. At what level is this wrong? If partial recovery is possible, what bounds it?

  \textit{Hint:} Recovery is another engine requiring a cold end, ultimately surroundings near \(300\ \mathrm{K}\). Even exhaust at \(700\ \mathrm{K}\) permits only \(1-300/700\approx 57\%\) Carnot recovery. Complete conversion with no colder heat rejection directly violates Kelvin.
\end{puzzles}

From an impossible ship we found thermal traffic rules: conservation governs quantity, the second law direction, Carnot bounds engine appetite with two temperatures, and the third law seals temperature's floor. The three thermal laws are now complete: zeroth establishes temperature, first supplies the ledger, second direction, third boundaries. Next we exchange macroscopic engines and reservoirs for molecular armies, exploring how pressure, temperature, and entropy emerge statistically. Chapter 29's microstate counting becomes a complete statistical physics.
